\documentclass[10pt,a4paper]{amsart}
\usepackage[margin=19mm]{geometry}
\usepackage{amsmath,amssymb,mathtools}
\usepackage[scr=rsfs,cal=euler]{mathalfa}
\usepackage{xfrac}
\usepackage[unicode,hidelinks,bookmarks=false]{hyperref}
\newcommand{\ie}{i.e.\ }
\newcommand{\cf}{cf.\ }
\newcommand{\resp}{respectively\ }
\newcommand{\Subsection}{\subsection}
\providecommand{\nobreakdash}{\nobreak\hbox{-}}
\setlength{\emergencystretch}{2em}
\multlinegap0pt
\usepackage{bm}
\usepackage{bbm}
\usepackage{dsfont}
\usepackage{xspace}
\usepackage{cancel}
\usepackage{mathtools}
\usepackage{amsthm}
\makeatletter
\@ifundefined{thm}{\newtheorem{thm}{Thm}}{}
\@ifundefined{prop}{\newtheorem{prop}[thm]{Proposition}}{}
\makeatother
\makeatletter
\expandafter\ifx\csname defn\endcsname\relax
\newenvironment{defn}{\refstepcounter{thm}\begin{proof}[Definition \emph{\thethm}]\renewcommand*{\qedsymbol}{\(\blacksquare\)}}{\end{proof}}
\fi
\expandafter\ifx\csname ex\endcsname\relax
\newenvironment{ex}{\refstepcounter{thm}\begin{proof}[Example \emph{\thethm}]\renewcommand*{\qedsymbol}{\(\blacksquare\)}}{\end{proof}}
\fi
\expandafter\ifx\csname quest\endcsname\relax
\newenvironment{quest}{\refstepcounter{thm}\begin{proof}[Question \emph{\thethm}]\renewcommand*{\qedsymbol}{\(\blacksquare\)}}{\end{proof}}
\fi
\expandafter\ifx\csname rem\endcsname\relax
\newenvironment{rem}{\refstepcounter{thm}\begin{proof}[Remark \emph{\thethm}]\renewcommand*{\qedsymbol}{\(\blacksquare\)}}{\end{proof}}
\fi
\expandafter\ifx\csname rethm\endcsname\relax
\newenvironment{rethm}[1]
  {\renewcommand\theinnercustomthm{#1}\innercustomthm}
  {\endinnercustomthm}
\fi
\expandafter\ifx\csname warn\endcsname\relax
\newenvironment{warn}{\refstepcounter{thm}\begin{proof}[Warning \emph{\thethm}]\renewcommand*{\qedsymbol}{\(\blacksquare\)}}{\end{proof}}
\fi
\makeatother
\title[Separating dots with circles]{Separating dots with circles}
\author{James Beyer, Jaewon Min, Greg Muller}
\date{Source: arXiv:2505.22851v1 (CC BY 4.0)}
\begin{document}
\maketitle

\noindent\emph{Source and licence.} Separating dots with circles. Version arXiv:2505.22851v1 (CC BY 4.0), distributed under the Creative Commons Attribution 4.0 International licence, as recorded in the arXiv licence metadata for this exact version and shown on the abstract page https://arxiv.org/abs/2505.22851v1. Full rights notice: licenses/arxiv-2505.22851-CC-BY-4.0.txt.\par\smallskip

\noindent\textit{Editorial scope.} Counted unit: the Prop whose source label is \texttt{prop: orientedincident} in the article (source0.tex lines 664--668, source label \texttt{prop: orientedincident}), delivered together with the article's own complete original proof of it (source0.tex lines 670--706, one \texttt{proof} environment of 37 lines). No mathematics is written, repaired or re-derived: statement and proof are the article's own text line for line. Required context retained uncounted: prop: incident0 (lines 636--638). Every definition, display and label of those lines is retained unchanged. Omitted from this package and not redistributed: 1--635, 639--663, 669--669, 707--1894. Nothing inside a retained excerpt is omitted or altered: every retained body is delivered exactly as the article's lines are written, so no omission receipt of any kind accompanies this package. Citations: the retained excerpts carry no citation command at all, so no bibliography entry of the article is transcribed and no reference is redistributed. Reference adaptation: none was needed and none was made; every reference inside the delivered unit resolves inside it, so no unresolved reference or citation is produced. Printed slips kept literal: no printed word, symbol, hypothesis or number of the counted statement or of its proof was altered. Layout and class: the article's own class is \texttt{amsart} with packages including amsmath, amsfonts, amsthm, amssymb, MnSymbol, tikz, setspace, fullpage, hyperref, subcaption; the wrapper is the lane's pinned \texttt{amsart} editorial wrapper at 10pt on A4, which restates the theorem-like environments the retained text needs and adds the article's own macro definitions that the retained text needs (none), with extra wrapper packages (bm, bbm, dsfont, xspace, cancel, mathtools). The delivered numbering is the unit's own. Assets: no retained span contains a figure environment, a graphics-inclusion command, a PDF-page inclusion, a table or a TikZ picture; no image file is copied into this package, the package declares no asset dependency and no image file is redistributed with it. Licence: version arXiv:2505.22851v1 of this article is distributed under the Creative Commons Attribution 4.0 International licence, as recorded in the arXiv licence metadata for this exact version and shown on the abstract page https://arxiv.org/abs/2505.22851v1. A full rights notice is delivered at licenses/arxiv-2505.22851-CC-BY-4.0.txt. Characters: every retained excerpt is pure ASCII, so the delivered engine, XeLaTeX with its default Latin Modern text font, reports no missing character.
\begin{prop}\label{prop: incident0}
Let $n\geq4$. Given $n$-many dots $D$ in general position in the sphere $S$, there are $(2n-4)$-many incident circles with no dots on one side.
\end{prop}

\begin{prop}\label{prop: orientedincident}
Let $n\geq3$ and $0\leq k \leq n-3$. Given $n$-many dots $D$ in general position in the sphere $S$, there are 
\[I_{k,n} := 2(k+1)(n-k-2)\]
oriented incident circles with $k$-many dots on their left side (and thus $n-k-3$ dots on the right).
\end{prop}

\begin{proof}
If $n=3$, then $k=0$. Since the unique incident circle through the three dots has $2=I_{0,3}$ orientations, the theorem holds in this case. %In what follows, assume $n\geq4$.% and so Proposition \ref{prop: incident0} applies.

We prove the $n\geq4$ case by induction on $k$. 
If $k=0$, then Proposition \ref{prop: incident0} implies there are $(2n-4)$-many incident circles with no dots on one side. Since $n\geq4$, the other side must have at least 1 dot and there is a unique orientation of each such circle which has no dots on the left side. Thus, there are $I_{0,n}=2\cdot 1\cdot(n-2)$-many oriented incident circles with $0$ dots on their left side.

Next, assume the proposition holds for $k-1$, and consider a configuration $D$ of $n$-many dots in general position in $S$. 
There are $n$-many ways to choose a dot $d\in D$ to delete; by assumption, there are then $I_{k-1,n-1}$-many ways to choose an oriented incident circle through three of the remaining $(n-1)$-many dots which has $(k-1)$-many dots on its left side. 
Therefore, there are
\[
n \cdot I_{k-1,n-1} = 2nk(n-k-2)
\]
ways to delete a dot and then choose an oriented incident circle with $k-1$ dots on its left side.

Alternatively, we could first choose the circle, then choose the dot to delete. This can be done in two disjoint ways.
    \begin{itemize}
        \item First choose an oriented incident circle with $(k-1)$-many dots on its left side, then choose a dot on that circle's right side to delete. By assumption, the number of ways to do this is
        \[
        I_{k-1,n}\cdot (n-k-2) = 2k(n-k-1)(n-k-2)
        \]
        \item First choose an oriented incident circle with $k$-many dots on its left side, then choose a dot on that circle's left side to delete. The number of ways to do this is
        \[
        (\text{$\#$ of oriented incident circles of $D$ with $k$-many dots on their left side})\cdot k
        \]
    \end{itemize}
Since this constructs the same set as before, the sum of these two counts is equal to $nI_{k-1,n-1}$. Solving this equality for the desired number of circles,
\[
% 2nk(n-k-3)
% =
% 2k(n-k-2)(n-k-3)
% +
(\text{$\#$ of oriented incident circles of $D$ with $k$-many dots on their left side})
%= 2n(n-k-3) - 2(n-k-1)(n-k-3)
=2(k+1)(n-k-2)%=:I_{k,n}
\]
This completes the induction.
\end{proof}

\end{document}
